\documentclass[11pt]{article}

\usepackage{lmodern}
\usepackage[T1]{fontenc}
\usepackage[margin=1in]{geometry}
\usepackage{amsmath}
\usepackage{amssymb}
\usepackage{bm}
\usepackage{hyperref}
\usepackage{enumitem}
\usepackage{booktabs}

\hypersetup{colorlinks=true, linkcolor=blue, urlcolor=blue, citecolor=blue, bookmarksopen=true}

\newcommand{\dd}{\mathrm{d}}
\newcommand{\pd}[2]{\frac{\partial #1}{\partial #2}}
\newcommand{\vecb}[1]{\bm{#1}}

\title{AeroHexMesh: Mathematical Formulations\\[0.3em]
\large A companion to \texttt{docs/} -- every algorithm's math, in one document}
\author{}
\date{}

\begin{document}
\maketitle
\tableofcontents

\section*{Scope}
\addcontentsline{toc}{section}{Scope}

This document collects, in standard mathematical notation, every
algorithm documented (in prose + pseudocode) across \texttt{docs/}. It
is organized identically to that tree: one section per pipeline
directory, in the same order. Each section cross-references the
corresponding Markdown page, which additionally covers \emph{where}
each piece of math lives in code and \emph{why} each design decision
was made -- this document is the math alone, typeset properly.

Notation used throughout: bold lowercase (\(\vecb{p}\), \(\vecb{n}\))
for 2D or 3D vectors/points; \(s\) for an arc-length parameter;
subscripts for grid/element indices; primes or \(\partial\) for
derivatives as is conventional per context.

%=====================================================================
\section{\texttt{Construct2D/}}
\label{sec:construct2d}
Construct2D's own grid-generation PDE (elliptic or hyperbolic,
selected by the \texttt{slvr} namelist field) is vendored upstream
Fortran and out of scope here -- see
\texttt{Construct2D/doc/user\_manual.pdf}. No original math from this
repository's own code is introduced in this directory; see
\href{https://github.com/kvishalrai/AeroHexMesh/blob/master/docs/Construct2D/README.md}{\texttt{docs/Construct2D/README.md}}
for the batch-driving mechanism.

%=====================================================================
\section{\texttt{Construct2D\_to\_SOD2D/linear\_mesh/}}
\label{sec:c2s-linear}
See \href{https://github.com/kvishalrai/AeroHexMesh/blob/master/docs/Construct2D_to_SOD2D/linear_mesh.md}{\texttt{docs/Construct2D\_to\_SOD2D/linear\_mesh.md}}.

\subsection{Spanwise extrusion}
Given a 2D structured grid \((x_{ij}, y_{ij})\), \(i \in [0,\text{idim}), j\in[0,\text{jdim})\),
and spanwise stations \(z_k = k \cdot \Delta z\), \(k \in [0, \text{nplanes})\),
extrusion is pure translation:
\begin{equation}
X_{ijk} = x_{ij}, \qquad Y_{ijk} = y_{ij}, \qquad Z_{ijk} = z_k .
\end{equation}
No offset or curvature is introduced -- contrast with
Section~\ref{sec:etagrid-sweep}'s curved sweep.

\subsection{Structured node numbering}
\label{sec:c2s-numbering}
For block \(n\) with per-block node count \(N_n = \text{idim}_n \cdot \text{jdim}_n \cdot \text{kdim}_n\),
the flat (1-based) node id of structured index \((i,j,k)\) is
\begin{equation}
\mathrm{id}(n,i,j,k) = \underbrace{1 + \sum_{m<n} N_m}_{\text{base}_n} \;+\; k + d_k j + d_k d_j i,
\end{equation}
where \(d_j, d_k\) are block \(n\)'s own J/K extents (K fastest-varying,
Fortran-style). Hexahedra are built by walking every interior
\((i,j,k)\) cell and gathering its 8 corners via a fixed
\((-1,0)\)-offset shift table in each axis.

\paragraph{O-grid periodic closure.} The last \(i\)-plane
(\(i=\text{idim}-1\)) is never stored; any reference to it is rewritten
as \(i \gets 0\) before evaluating \(\mathrm{id}(\cdot)\) -- an
\emph{aliasing} closure, not a duplicated node with a periodic link.

\paragraph{C-grid wake-cut fold.} A sub-range of the wall row
(\(j=0\)) is aliased onto another sub-range of the \emph{same} row via
a constant index offset:
\begin{equation}
i_B \;=\; \omega - i_A, \qquad \omega \;=\; s_{2,B} - s_{2,A},
\end{equation}
where \(s_{2,A}, s_{2,B}\) are the two matched NMF sub-range start
indices. Side \(B\)'s nodes are never stored; every reference to an
index in \([i_{\mathrm{lo},B}, i_{\mathrm{hi},B}]\) at \(j=0\) is
rewritten through this map first.

\subsection{Inlet/outlet classification by angle of attack}
\label{sec:aoa-classify}
At each farfield boundary point \((i,k)\), with wall-row point
\(\vecb{p}_w = (x_{i,0,k}, y_{i,0,k})\) and farfield-row point
\(\vecb{p}_f = (x_{i,\text{jmax},k}, y_{i,\text{jmax},k})\):
\begin{align}
\hat{\vecb{d}} &= \frac{\vecb{p}_f - \vecb{p}_w}{\lVert \vecb{p}_f - \vecb{p}_w \rVert}
&&\text{(local outward direction)} \\
\vecb{U}_\infty &= (\cos\alpha,\ \sin\alpha)
&&\text{(free-stream direction, } \alpha = \text{AoA, mesh built at } \alpha=0\text{)} \\
\text{class} &=
\begin{cases}
\text{inlet}, & \hat{\vecb{d}} \cdot \vecb{U}_\infty < 0 \quad \text{(flow enters)}\\
\text{outlet}, & \hat{\vecb{d}} \cdot \vecb{U}_\infty \ge 0 \quad \text{(flow leaves)}
\end{cases}
\end{align}
For a C-grid, the two wake-end faces are always classified outlet,
bypassing this test entirely (their own outward direction is not
meaningfully wall\(\to\)farfield). \textbf{This sign convention has not
been independently verified} -- see
\href{https://github.com/kvishalrai/AeroHexMesh/blob/master/docs/reference/troubleshooting.md}{\texttt{docs/reference/troubleshooting.md}}.

\subsection{Wall-spline fit: natural cubic spline}
\label{sec:wall-spline}
Given ordered wall corner points \((x_i, y_i)\), the arc-length
parametrization is cumulative chord length,
\begin{equation}
s_0 = 0, \qquad s_i = s_{i-1} + \lVert \vecb{p}_i - \vecb{p}_{i-1} \rVert .
\end{equation}
An open natural cubic spline is fit to \(y(s)\) (and independently to
\(x(s)\)); on segment \(i\), for \(s_i \le t \le s_{i+1}\), \(\delta t = t - s_i\):
\begin{equation}
y(t) = a_i + b_i\,\delta t + c_i\,\delta t^2 + d_i\,\delta t^3 .
\end{equation}
The second-derivative values \(M_i = y''(s_i)\) solve the classic
tridiagonal system with natural boundary conditions \(M_0 = M_n = 0\):
\begin{equation}
h_i = s_{i+1}-s_i, \qquad
h_{i-1} M_{i-1} + 2(h_{i-1}+h_i) M_i + h_i M_{i+1}
= 6\!\left[\frac{y_{i+1}-y_i}{h_i} - \frac{y_i - y_{i-1}}{h_{i-1}}\right],
\end{equation}
solved by the Thomas algorithm (tridiagonal Gaussian elimination:
forward sweep then back substitution). The per-segment coefficients
follow directly:
\begin{align}
a_i &= y_i, &
b_i &= \frac{y_{i+1}-y_i}{h_i} - \frac{h_i(2M_i + M_{i+1})}{6}, \\
c_i &= \frac{M_i}{2}, &
d_i &= \frac{M_{i+1}-M_i}{6h_i}.
\end{align}

\subsection{Gmsh periodicity}
The spanwise periodic pair is declared as a pure translation:
\begin{equation}
\text{Periodic Surface } \{k_{\max}\} = \{k_0\}\ \text{Translate}\ (0,0,\text{span}_z).
\end{equation}

%=====================================================================
\section{\texttt{Construct2D\_to\_SOD2D/high\_order\_mesh/}}
\label{sec:c2s-high-order}
See \href{https://github.com/kvishalrai/AeroHexMesh/blob/master/docs/Construct2D_to_SOD2D/high_order_mesh.md}{\texttt{docs/Construct2D\_to\_SOD2D/high\_order\_mesh.md}}.

\subsection{Parametric (arc-length) high-order wall placement}
\label{sec:parametric-placement}
For a high-order wall node with reference (Gauss--Lobatto--Legendre)
coordinate \(\xi_k \in [-1,1]\) along a face's own along-wall axis,
between two corners matched to spline arc-lengths \(s_{\mathrm{lo}}\),
\(s_{\mathrm{hi}}\) (segment index \(\mathrm{iseg} = \min\)):
\begin{align}
s_f &= s_{\mathrm{lo}} + (s_{\mathrm{hi}} - s_{\mathrm{lo}})\cdot\tfrac{1}{2}(\xi_k + 1), \\
\delta t &= s_f - s_{\mathrm{iseg}}, \\
x_t &= a_x + \delta t\,(b_x + \delta t\,(c_x + \delta t\, d_x)), \\
y_t &= a_y + \delta t\,(b_y + \delta t\,(c_y + \delta t\, d_y)), \\
\vecb{u}_{\text{buffer}} &= \lambda_{\text{scale}} \cdot \big((x_t,y_t) - (x_{\text{current}}, y_{\text{current}})\big),
\end{align}
where \((a,b,c,d)_{x,y}\) are the segment's cubic coefficients from
Section~\ref{sec:wall-spline} and \(\lambda_{\text{scale}}\) is an optional
damping factor (default 1, no damping). This is a
\emph{parametric}, not nearest-point, placement -- deliberately, since
a Euclidean nearest-point search is a many-to-one map near high
curvature.

\subsection{Linear elasticity relaxation}
\label{sec:elasticity-pde}
The imposed wall displacement becomes a Dirichlet condition; the
interior solves the isotropic linear elasticity equilibrium equation
\begin{equation}
\nabla\cdot\vecb{\sigma} = \vecb{0}, \qquad
\vecb{\sigma} = \lambda(\nabla\cdot\vecb{u})\,\vecb{I} + 2\mu\,\vecb{\varepsilon}(\vecb{u}), \qquad
\vecb{\varepsilon}(\vecb{u}) = \tfrac{1}{2}\!\left(\nabla\vecb{u} + \nabla\vecb{u}^{\!\top}\right),
\end{equation}
with the Lam\'e parameters computed directly from Young's modulus
\(E\) and Poisson's ratio \(\nu\):
\begin{equation}
\mu = \frac{E}{2(1+\nu)}, \qquad
\lambda = \frac{E\nu}{(1+\nu)(1-2\nu)} .
\end{equation}
Discretized with the same spectral-element basis the flow solver
itself uses, solved by preconditioned conjugate gradient (block-Jacobi
or scalar-Laplacian-diagonal preconditioner). After solving, if the
resulting quality \(\min Q < 10^{-6}\) (an inverted/degenerate
element), the solver reverts and performs a small grid search over
\((E,\nu)\) (coarse 2-point, then finer 9-point), re-solving each time
and keeping whichever combination maximizes \(\min Q\).

\subsection{Opposite-face-distance displacement (\texttt{opposite\_face\_disp\_factor}, in progress)}
\label{sec:opposite-face}
A second, independent displacement mode: each wall node is displaced by
an integer multiple of the distance to its own \emph{opposite-face
node} -- the node directly across the same hex element, found by
\texttt{mod\_wall\_model.f90}'s \texttt{evalEXAtFace} (the same routine
SOD2D's LES wall models use for exchange-distance evaluation, reused
verbatim, not reimplemented):
\begin{equation}
\vecb{u}_{\text{buffer}} = \lambda_{\text{scale}} \cdot n_{\text{factor}} \cdot L \cdot \hat{\vecb{n}},
\end{equation}
where \(L\) is the through-element distance to the opposite-face node,
\(\hat{\vecb{n}}\) its inward (into-the-fluid) normal, \(n_{\text{factor}}
= \texttt{opposite\_face\_disp\_factor}\) (integer), and
\(\lambda_{\text{scale}}\) reuses the same substep-damping factor as
Section~\ref{sec:parametric-placement}'s parametric placement.

\paragraph{Known limitation (not yet usable at \(n_{\text{factor}}=1\)).}
The automatic \((E,\nu)\) fallback search above exhausts without finding
\(\min Q \geq 10^{-6}\) on every mesh tested (a periodic NACA0012
C-grid-like mesh, a Construct2D O-grid mesh, and a fully 3D M6 wing mesh
with no periodicity and a blunt trailing edge) -- not specific to one
mesh topology.

\paragraph{Diagnostic sequence.} Four candidate mechanisms were tested
directly against dumped data, in order, before the actual cause was
found:
\begin{enumerate}
\item \emph{Simple normal-gap closure.} Refuted: dumping each wall
  node's own position alongside its opposite-face node's position,
  before and after the elasticity solve, shows the interior node
  follows the wall almost exactly (median displacement-magnitude ratio
  \(0.998\), median direction cosine \(1.0000\)); the gap barely
  changes for most nodes (median \(1.7128\times10^{-3} \to
  1.7092\times10^{-3}\)).
\item \emph{In-plane (streamwise) shear} between neighboring wall
  stations. Refuted: the signed area of the quadrilateral
  (wall\(_1\), wall\(_2\), opp\(_2\), opp\(_1\)) for every nearest-
  neighbor wall-node pair, before vs.\ after -- 0 of 6604 sampled quads
  flipped sign or went near-zero; worst shrink ratio \(0.57\).
\item \emph{Spanwise (Z) variation} of the imposed field. Refuted, and
  confirmed \emph{correct} instead: magnitude spread across z-stations
  at matching \((x,y)\) was exactly \(0.0000\times10^{0}\), direction
  spread \(3.3\times10^{-16}\) -- perfectly z-invariant, as required.
\item Only then: directly locating the failing element(s), by dumping
  \texttt{computeQuality}'s per-element \texttt{quality(:)} and
  per-Gauss-point \texttt{quality\_pv} (reused via \texttt{mu\_e(:,:)})
  for elements below a threshold.
\end{enumerate}

\paragraph{The two failure signatures found.}
\begin{itemize}
\item \textbf{Trailing-edge tangent degeneracy} (NACA0012 meshes):
  exactly \textbf{12 elements}, all with quality and \texttt{mu\_e}
  exactly \(0\), all wall-touching, clustered at \(x\approx1.0,
  y\approx0\) (the TE, upper and lower surface, every spanwise
  station). The imposed displacement at the TE tip node itself:
  \((d_x,d_y,d_z)\approx(+1.71\times10^{-3}, -1.4\times10^{-7}, 0)\) --
  almost \emph{purely tangential} (downstream), not into the fluid at
  all. A sharp/near-sharp TE has no well-defined surface normal, so
  \(\hat{\vecb{n}}\) degenerates toward the local element's chordwise
  axis there instead of a true perpendicular, shearing an already
  acutely-angled element along its own edge until its Jacobian
  inverts, regardless of \((E,\nu)\) -- a \emph{direction} problem, not
  a magnitude one.
\item \textbf{Wing-root band collapse} (M6 wing mesh, no periodicity,
  blunt TE -- a genuinely different case): \textbf{88} exactly-zero-
  quality elements, all in the single row of wall elements closest to
  the root symmetry plane (\(z\approx0.0091\)--\(0.0101\)), spanning
  almost the \emph{entire chord} (\(x=0.005\) to \(x\approx0.81\)), not
  one corner. The individual displacement vectors there look
  geometrically sane in isolation (magnitude \(\sim7\times10^{-5}\),
  \(d_z\) exactly \(0\), direction transitioning smoothly from
  \(-x\)-dominant near the LE to \(y\)-dominant mid-chord) -- a locally
  sane vector field that still collapses the whole row it acts on.
  \textbf{Not yet root-caused.}
\end{itemize}
Substepping (smaller \(\lambda_{\text{scale}}\), more substeps) does
not fix either failure -- confirmed by running the fallback search on a
single substep at \(\lambda_{\text{scale}}=0.3\), which still fails
identically. Expected for the TE case: a magnitude reduction cannot fix
a displacement pointed the wrong way.

%=====================================================================
\section{\texttt{Construct2D\_to\_NEKRS/linear\_mesh/}}
\label{sec:c2n-linear}
See \href{https://github.com/kvishalrai/AeroHexMesh/blob/master/docs/Construct2D_to_NEKRS/linear_mesh.md}{\texttt{docs/Construct2D\_to\_NEKRS/linear\_mesh.md}}.
Extrusion/NMF-remap/node-numbering/AoA-classification math is shared
with Section~\ref{sec:c2s-linear}--\ref{sec:aoa-classify}. This section
covers what is genuinely different: Nek5000's own conversion chain and
Fischer's independently-implemented smoothing algorithm.

\subsection{Format-conversion chain (\texttt{arglist.sh})}
The spanwise extrusion itself happens \emph{inside} Nek5000's own
tools, not in Python:
\begin{equation}
\texttt{.re2} \xrightarrow{\texttt{re2torea}} \texttt{.rea (2D)}
\xrightarrow{\texttt{n2to3}} \texttt{.rea (3D, extruded+periodic)}
\xrightarrow{\texttt{reatore2}} \texttt{.re2 (3D)}.
\end{equation}

\subsection{Case-parameter extraction from the NMF}
Reading the native 2D NMF's own \texttt{VISCOUS} entry
(\(f_1=3\), excluding the C-grid wake-cut \texttt{ONE\_TO\_ONE} pairing):
\begin{align}
\text{ring\_width} &= \text{idim} - 1, &
\text{radial\_rings} &= \text{jmax} - 1, \\
\text{wall\_start\_element} &= s_1^{\text{VISCOUS}}, &
\text{wall\_element\_count} &= e_1^{\text{VISCOUS}} - s_1^{\text{VISCOUS}} .
\end{align}

\subsection{Fischer's natural cubic spline (\texttt{spline}/\texttt{splint})}
\label{sec:fischer-spline}
Independently implemented (Fortran77, Numerical Recipes style) but
mathematically the \emph{same} construction as Section~\ref{sec:wall-spline}:
natural boundary conditions, tridiagonal solve for \(y''\). Evaluation
(\texttt{splint}) uses the equivalent cubic-Hermite-style blend, given
bracketing points \(x_{\mathrm{lo}}, x_{\mathrm{hi}}\) (\(h = x_{\mathrm{hi}}-x_{\mathrm{lo}}\)):
\begin{align}
a &= \frac{x_{\mathrm{hi}} - x}{h}, \qquad b = \frac{x - x_{\mathrm{lo}}}{h}, \\
y &= a\,y_{\mathrm{lo}} + b\,y_{\mathrm{hi}} + \big[(a^3-a)\,y''_{\mathrm{lo}} + (b^3-b)\,y''_{\mathrm{hi}}\big]\frac{h^2}{6}.
\end{align}
Every high-order node is placed at its own GLL-reference arc-length
fraction along its segment,
\begin{equation}
s_f(j,e) = s_{i-1} + \tfrac{1}{2}\,\mathrm{ds}\,(\zeta_j + 1), \qquad \zeta_j = \text{the GLL reference coordinate} \in[-1,1],
\end{equation}
the same conceptual move as Section~\ref{sec:parametric-placement} --
independently arrived at in a separate codebase.

\subsection{Weighted-Laplace boundary blending}
\label{sec:weighted-laplace}
Given the wall's own displacement field \(\delta u\) (from
Section~\ref{sec:fischer-spline}), extend it into the volume by solving a
\emph{weighted Laplace} equation (not the full elasticity tensor of
Section~\ref{sec:elasticity-pde}):
\begin{equation}
\nabla\cdot\big(h_1(\vecb{x})\,\nabla u\big) = 0, \qquad u = \delta u \ \text{on the wall}.
\end{equation}
Boundary-layer protection via a spatially-varying diffusion coefficient:
\begin{align}
\delta &= \frac{\text{volume(wall-adjacent elements)}}{\text{surface area(wall-adjacent elements)}}
&&\text{(avg. boundary-layer element thickness)}\\
\delta_p &= 5\delta && \text{(``protected'' thickness)} \\
h_1(\vecb{x}) &= 1 + 9\exp\!\left(-\Big(\frac{d(\vecb{x})}{\delta_p}\Big)^{\!2}\right)
&&\text{(} d = \text{distance to wall; } h_1 \to 10 \text{ at the wall, } \to 1 \text{ far away)}
\end{align}
Solved via the standard Dirichlet-lift decomposition \(u = u_0 + u_b\)
(\(u_b\) = boundary data extended by duplicate-node averaging;
\(u_0=0\) on Dirichlet boundaries):
\begin{equation}
\bar A\, u_0 = -M \bar A\, u_b,
\end{equation}
where \(\bar A\) is the (weighted) Laplace/Helmholtz operator and
\(M\) the interior/Dirichlet mask; \(x\)- and \(y\)-displacement are
solved as two independent scalar problems (no shear/Poisson coupling
between components, unlike Section~\ref{sec:elasticity-pde}).

%=====================================================================
\section{\texttt{Construct2D\_to\_NEKRS/high\_order\_mesh/}}
\label{sec:c2n-high-order}
See \href{https://github.com/kvishalrai/AeroHexMesh/blob/master/docs/Construct2D_to_NEKRS/high_order_mesh.md}{\texttt{docs/Construct2D\_to\_NEKRS/high\_order\_mesh.md}}.
No new smoothing math: this step \emph{interpolates} Section~\ref{sec:c2n-linear}'s
already-corrected, fixed-order (\(\ell x_1=4\)) reference sampling onto
the production mesh's own (possibly different) polynomial order via a
standard Lagrange-basis change-of-interpolant operation
(\texttt{map\_m\_to\_n}: an \(m\)-point to \(n\)-point GLL resampling),
then copies that corrected \((x,y)\) identically to every spanwise
plane (pure-translation extrusion, so no recomputation needed there):
\begin{equation}
x_{m1}(i,j,i_z,e) = x_{m1}(i,j,1,e) \quad \forall\, i_z = 2,\dots,\ell z_1 .
\end{equation}

%=====================================================================
\section{\texttt{pyHyp/}}
\label{sec:pyhyp}
See \href{https://github.com/kvishalrai/AeroHexMesh/blob/master/docs/pyHyp/README.md}{\texttt{docs/pyHyp/README.md}}.
pyHyp's own hyperbolic marching PDE is an external tool, out of scope.
This section covers this repo's own surface-preprocessing math.

\subsection{Surface coarsening}
Keep every other structured-grid index, plus the last index
unconditionally (so both endpoints of an axis are always retained):
\begin{equation}
\text{keep}(n) = \{0, 2, 4, \dots\} \cup \{n-1\}.
\end{equation}

\subsection{Local surface curvature}
\label{sec:curvature}
Given a per-block bicubic Hermite surface fit \(\vecb{P}(s,t)\)
(evaluated per-cell, finite-difference step \(\varepsilon\)):
\begin{align}
\vecb{P}_s &= \pd{\vecb{P}}{s}, \quad \vecb{P}_t = \pd{\vecb{P}}{t}
&&\text{(central difference)} \\
\vecb{P}_{ss}, \vecb{P}_{tt}, \vecb{P}_{st} &&&\text{(central-difference 2nd partials)} \\
\hat{\vecb{n}} &= \frac{\vecb{P}_s \times \vecb{P}_t}{\lVert \vecb{P}_s \times \vecb{P}_t \rVert}
&&\text{(unit normal)}
\end{align}
First fundamental form \(E,F,G\) and second fundamental form \(L,M,N\):
\begin{equation}
E = \vecb{P}_s\!\cdot\!\vecb{P}_s,\ \ F = \vecb{P}_s\!\cdot\!\vecb{P}_t,\ \ G = \vecb{P}_t\!\cdot\!\vecb{P}_t;
\qquad
L = \vecb{P}_{ss}\!\cdot\!\hat{\vecb{n}},\ \ M = \vecb{P}_{st}\!\cdot\!\hat{\vecb{n}},\ \ N = \vecb{P}_{tt}\!\cdot\!\hat{\vecb{n}} .
\end{equation}
Mean curvature \(H\), Gaussian curvature \(K\), and total curvature
magnitude \(\kappa\):
\begin{equation}
H = \frac{EN + GL - 2FM}{2(EG-F^2)}, \qquad
K = \frac{LN - M^2}{EG - F^2}, \qquad
\kappa = \sqrt{\max(4H^2 - 2K,\ 0)} \;\big(= \sqrt{k_1^2+k_2^2}\big).
\end{equation}

\subsection{Curvature monitor field}
Per-cell curvature is clipped at its own 99.5th percentile (outlier
robustness), normalized by the block's own median \(\kappa_{\mathrm{ref}}\),
scaled by an aggressiveness parameter \(\beta\), and capped:
\begin{equation}
m_{\text{cell}} = \min\!\left(1 + \beta\,\frac{\kappa_{\text{cell}}}{\kappa_{\mathrm{ref}}},\ \ m_{\max}\right),
\end{equation}
then averaged onto nodes over each node's up-to-4 surrounding cells.
\(\beta=0 \Rightarrow m \equiv 1\) (recovers plain, uniform elliptic
smoothing).

\subsection{Curvature-weighted 1D equidistribution}
\label{sec:equidistribution}
Redistributes a curve's own points to equal increments of
\emph{weighted} arc length. With chord-length parameter \(s\), total
length \(S\), and monitor \(m(s)\) (linearly interpolated from a fine
resampling):
\begin{equation}
W(s) = \int_0^s m(\sigma)\,\dd\sigma, \qquad
s_i = W^{-1}\!\left(\frac{i}{n-1}\,W(S)\right), \quad i=0,\dots,n-1,
\end{equation}
monotonic by construction; \(m\equiv 1\) reproduces uniform arc-length
spacing exactly.

\subsection{Weighted elliptic relaxation + Newton reprojection}
\label{sec:newton-reprojection}
Each interior point is damped-relaxed toward its curvature-weighted
4-neighbor average (the discrete equilibrium of a spring network with
local stiffness set by curvature), then reprojected onto the true
surface via local Gauss--Newton refinement of its own \((s,t)\)
location:
\begin{align}
\vecb{r} &= \vecb{x}_{\text{target}} - \vecb{P}(s,t), \\
\vecb{J} &= \big[\vecb{P}_s,\ \vecb{P}_t\big] \quad\text{(a } 3\times2 \text{ matrix)}
&&\text{(central-difference Jacobian)} \\
\delta &= \big(\vecb{J}^{\!\top}\vecb{J} + \eta \vecb{I}\big)^{-1} \vecb{J}^{\!\top}\vecb{r},
\qquad \eta = 10^{-12}
&&\text{(regularized normal equations)} \\
(s,t) &\leftarrow (s,t) + \delta .
\end{align}
\(s,t\) are allowed to run outside \([0,1]\) mid-solve; an overflow
converts into an exact cell-index shift (cells are \(C^0\)-continuous
at \(s=1\)/\(t=1\)), and the Newton solve resumes in the new cell.

%=====================================================================
\section{\texttt{etaGrid\_to\_SOD2D/}}
\label{sec:etagrid}
See \href{https://github.com/kvishalrai/AeroHexMesh/blob/master/docs/etaGrid_to_SOD2D/README.md}{\texttt{docs/etaGrid\_to\_SOD2D/README.md}}.
The underlying \(\eta\)-grid concept -- wall-normal/spanwise grid sizes
proportional to the local Kolmogorov scale \(\eta\) -- is from Rouhi,
Kumar, Wu, Kozul \& Lehmkuhl, ``Leveraging unstructured grids for direct
numerical simulations of wall turbulence,'' under review, \emph{Journal
of Fluid Mechanics}
(\href{https://arxiv.org/abs/2605.01015}{arXiv:2605.01015}); what
follows is this repository's own mesh-generation and SOD2D-import math
built around that concept.

\subsection{Airfoil curve: fit and differential geometry}
Arc-length parametrization as in Section~\ref{sec:wall-spline}; a full cubic
spline \(\vecb{P}(s) = (\mathrm{spline}_x(s), \mathrm{spline}_y(s))\)
(\texttt{scipy.interpolate.CubicSpline}, not-a-knot end conditions).
Tangent and outward normal:
\begin{equation}
\hat{\vecb{t}}(s) = \frac{\vecb{P}(s+\varepsilon) - \vecb{P}(s-\varepsilon)}{2\varepsilon}
\Big/ \Big\lVert \cdot \Big\rVert, \qquad
\hat{\vecb{n}}(s) = \big(t_y(s),\ -t_x(s)\big) \quad (\text{tangent rotated } -90^\circ).
\end{equation}

\subsection{Vinokur (1983) two-sided stretching}
\label{sec:vinokur}
Redistributes \(n\) points along a segment with independently
specified first/last spacing \(\delta s_0, \delta s_1\) (as fractions
of the segment's own total length; \(s_0 = \delta s_0 (n-1)\),
\(s_1 = \delta s_1(n-1)\)):
\begin{equation}
A = \sqrt{s_1/s_0}, \qquad B = \frac{1}{\sqrt{s_0 s_1}} .
\end{equation}
\textbf{Case \(B>1\)} (spacings smaller than uniform): solve
\(\sinh(\Delta y)/\Delta y = B\) for \(\Delta y\) by Newton's method
(seeded by a small-\((B-1)\) series expansion, or
\(\log(2B)+\log\log(2B)\) for large \(B\)); then, with
\(\xi = i/(n-1)\),
\begin{equation}
u(\xi) = \frac{1}{2}\left(1 + \frac{\tanh\!\big(\Delta y(\xi - \tfrac12)\big)}{\tanh(\Delta y/2)}\right).
\end{equation}
\textbf{Case \(B<1\)}: the same construction with \(\sin\)/\(\tan\) in
place of \(\sinh\)/\(\tanh\), solving \(\sin(\Delta y)/\Delta y = B\).
\textbf{Case \(B\approx1\)}: \(u(\xi) = \xi\) (uniform). Finally,
\begin{equation}
t = \frac{u}{A + (1-A)u}, \qquad t_0 = 0,\ t_{n-1} = 1 \ \text{(exact)} .
\end{equation}
Points are evaluated at \(t\cdot(\text{total length})\) through a
cubic-spline fit of the \emph{original} segment, so resampled points
still lie exactly on the input curve.

\subsection{Body-wrap sweep geometry}
\label{sec:etagrid-sweep}
Given a YZ cross-section local coordinate \((Y_j, Z_j)\) at arc-length
station \(s_k\):
\begin{align}
X(k,j) &= P_x(s_k) + Y_j\, n_x(s_k), \\
Y(k,j) &= P_y(s_k) + Y_j\, n_y(s_k), \\
Z(k,j) &= Z_j,
\end{align}
i.e. a ruled-surface-style offset sweep: at \(Y_j=0\) every station
lands exactly on the true curve; away from the wall, each node is
carried at a fixed offset \emph{along the local normal at that
station} -- contrast with Section~\ref{sec:c2s-linear}'s pure-translation
extrusion.

\paragraph{Orientation.} A swept hex's signed volume,
\begin{equation}
V = (\vecb{e}_1 \times \vecb{e}_2)\cdot \vecb{e}_3, \qquad
\vecb{e}_1 = \vecb{p}_1-\vecb{p}_0,\ \vecb{e}_2=\vecb{p}_3-\vecb{p}_0,\ \vecb{e}_3=\vecb{p}_4-\vecb{p}_0,
\end{equation}
can come out uniformly negative across an entire swept region
(handedness ambiguity between the cross-section's own 2D winding and
the local sweep direction); corrected by a node-order swap
\([4,5,6,7,0,1,2,3]\) when the majority of a sampled batch is negative.

\subsection{Wake growth: geometric offset series}
\label{sec:wake-growth}
Downstream station offsets \(x_0=0, x_1=s_0, x_2=s_0(1+r), \dots\)
reach \(\text{total}\) in exactly \(n\) layers, where the ratio \(r\)
solves
\begin{equation}
f(r) = s_0\,\frac{r^n - 1}{r-1} = \text{total} \qquad (f(r)\to s_0 n \text{ as } r\to1),
\end{equation}
via bisection over \(r\in[1,\infty)\) (no closed form for general
\(n\)).

\subsection{Boundary classification (local coordinates)}
\begin{equation}
\text{wall}: |Y-Y_{\min}|<\text{tol}, \quad
\text{outer}: |Y-Y_{\max}|<\text{tol}, \quad
\text{periodic}_{0,1}: |Z-Z_{\min,\max}|<\text{tol},
\end{equation}
applied identically at every swept station since it depends only on
the cross-section's own local \((Y,Z)\), never on 3D position.

\subsection{Wavy-wall perturbation}
\label{sec:wavy-wall}
A second, independent displacement applied on top of an already
spline-smoothed wall, along the local outward normal:
\begin{equation}
\Delta(s,z) = A_{\text{frac}}\cdot c \cdot \mathrm{tent}\!\left(n_w\frac{z-z_{\min}}{z_{\text{span}}}\right)\cdot \tau(s),
\end{equation}
where \(c\) is the chord (auto-computed as the spline table's own
\(X\)-extent), \(A_{\text{frac}}\) an amplitude fraction, \(n_w\) an
integer wavenumber.

\paragraph{Chordwise taper.}
\begin{equation}
\tau(s) = \sin^2\!\left(\pi \frac{s}{s_{\text{total}}}\right),
\end{equation}
zero \emph{with zero slope} at both \(s=0\) and \(s=s_{\text{total}}\)
(both ends of the arc-length parametrization are the trailing edge),
peaking at the leading edge \(s = s_{\text{total}}/2\).

\paragraph{Spanwise unipolar tent wave.} With phase
\(\varphi = n_w (z-z_{\min})/z_{\text{span}} \bmod 1\):
\begin{equation}
\mathrm{tent}(\varphi) = 1 - |2\varphi - 1|, \qquad \varphi\in[0,1),
\end{equation}
range \([0,1]\): zero (the \emph{original}, undisplaced surface -- the
trough) at \(\varphi=0,1\); peak \(1\) at \(\varphi=0.5\). Because
\(n_w\) is an integer, \(\varphi\) at \(z=z_{\max}\) is exactly an
integer, so the wave is \emph{exactly} zero at both spanwise-periodic
boundaries.

\paragraph{Local tangent (Fortran evaluation).} Reusing the cubic
segment coefficients of Section~\ref{sec:wall-spline}:
\begin{equation}
t_x = b_x + \delta t\,(2c_x + 3\delta t\, d_x), \qquad
t_y = b_y + \delta t\,(2c_y + 3\delta t\, d_y), \qquad
\hat{\vecb{n}} = \frac{(t_y,-t_x)}{\lVert(t_x,t_y)\rVert} .
\end{equation}
The full imposed displacement is then
\(\vecb{u}_{\text{buffer}} = \Delta(s,z)\,\hat{\vecb{n}}\), with
\(u_{\text{buffer},z} \equiv 0\) always (displacement confined to the
local \(X\)–\(Y\) plane at every \(Z\)).

\subsection{Wall-node arc-length recovery (verification)}
\label{sec:newton-search}
For an already-placed wall node, recover \(s\) via Newton/Gauss--Newton
search minimizing \(\lVert \vecb{P}(s) - \vecb{x}\rVert^2\):
\begin{equation}
f(s) = \big(\vecb{P}(s)-\vecb{x}\big)\cdot\hat{\vecb{t}}(s), \qquad
f'(s) \approx \hat{\vecb{t}}(s)\cdot\hat{\vecb{t}}(s) \quad\text{(Gauss--Newton)}, \qquad
s \leftarrow s - f(s)/f'(s),
\end{equation}
seeded from the nearest of the curve's own control points; converges
in a handful of steps given a point already close to the curve.

%=====================================================================
\section{Cross-cutting: boundary conditions}
\label{sec:bc-reference}
See \href{https://github.com/kvishalrai/AeroHexMesh/blob/master/docs/reference/boundary-conditions.md}{\texttt{docs/reference/boundary-conditions.md}}
for the full physical-id table and \texttt{bc\_type} string reference
(not reproduced here as it is not mathematical content); the
inlet/outlet classification math is Section~\ref{sec:aoa-classify} above,
shared verbatim by every Gmsh-based pipeline that needs it.

\end{document}
